\section{Kac--Moody 局部化}\srcpage{49}
\providecommand{\Loc}{\mathrm{Loc}}% 本节新增宏，已在 reports/s11.md 登记
\providecommand{\Poinc}{\mathrm{Poinc}}% 本节新增宏，已在 reports/s11.md 登记

有了上一个图中那项重要的保留，我们迄今所讲的故事在任意层论中大体上都平行地成立；在很大的程度上，它在各种类似的设定中也成立，例如在 Fargues--Fontaine 曲线上。然而，几何 Langlands 猜想的证明用到一项只在 D-模理论中才具备的关键输入：Kac--Moody 局部化。这些想法可以追溯到 Beilinson 与 Drinfeld（1996）的开创性工作，他们从 oper 的概念出发，给出了 $\Bun_G$ 上 Hecke 特征层的一个构造。这是目前所知构造（尖点）Hecke 特征层最通用的办法，也是从 Galois 侧到自守侧这一方向的关键。这里，一个 oper 指一个（de Rham）$\Gd$-局部系统，连同一个把其底层的 $\Gd$-丛约化到 $\Ud^{\Th}$ 的约化，使得联络在相应的 $\Td$-丛上具有典范的形状。例如，当 $\Gd=\mathrm{SL}_2$ 时，它指的是秩 $2$ 的丛 $E$ 的一个滤过
\[
0\to\Th\to E\to\Th^{-1}\to 0,
\]
使得联络诱导出恒等映射
\[
\Th\to\Th^{-1}\otimes\Om^{1}=\Th.
\]

由于这一部分也是我最感生疏的部分，我在这里的讲述会比在本文其余部分更为印象式。基本想法是得到下图的一个版本
\[
\begin{tikzcd}[column sep=large,row sep=large]
D_{\qc}(*/\Gd)_{\Ran} \arrow[r,"\cong"] \arrow[d] & \Whit(G) \arrow[d] \\
D_{\qc}(\LocSys_{\Gd}) \arrow[r] & D(\Bun_G),
\end{tikzcd}
\]
只不过要交换 $G$ 与 $\Gd$ 的角色。

在谱侧，我们有某个隐约让人想起 $\Whit(G)$ 的东西：\srcpage{50}

\begin{definition}
形式圆盘 $\Spec(k[[t]])$ 上无单值 opers $\Op^{\monfree}_{\Gd}$ 的叠，是形式圆盘上带平坦联络的 $\Gd$-丛的模空间，并连带去心形式圆盘 $\Spec(k((t)))$ 上的一个 oper 结构。
\end{definition}

注意，形式圆盘上带平坦联络的 $\Gd$-丛全是平凡的；更确切地说，它们等价地由中心点处的一个 $\Gd$-挠子给出。叠 $\Op^{\monfree}_{\Gd}$ 无论在 Ind 意义下还是在 Pro 意义下都是极度无限型的。在所期望的「基本局部等价」的一侧，我们将有某个范畴
\[
\IndCoh^{*}(\Op^{\monfree}_{\Gd}),
\]
即 $\Op^{\monfree}_{\Gd}$ 上的归纳凝聚层构成的范畴；注意，由于这个对象极度无限维，要定义正确的归纳凝聚层概念并非易事。

在 Langlands 对应的另一侧，有一类由 Kazhdan--Lusztig 研究过的 Kac--Moody 表示构成的范畴，因而记作 $\KL(G)$。再次忽略各种 $\rho$-平移，它大致如下给出。回忆在有限维情形，李代数 $\gc$ 的表示也可以描述为 $*/\hat{G}$ 上的拟凝聚层，其中 $\hat{G}$ 是由 $\gc$ 生成的形式代数群。于是把表示从李代数 $\gc$ 积分到群 $G$，对应于沿 $*/\hat{G}\to*/G$ 延拓拟凝聚层。Kac--Moody 李代数的表示则与分类空间 $*/\LG^{\wedge}$ 上的拟凝聚层相关，这里 $\LG^{\wedge}$ 是 $\LG$ 在原点处的形式完备化。我们将关心其中能积分成正圈群 $\LpG\subset\LG$ 的表示的那些。这些表示也可以理解为分类空间
\[
*/(\LpG\subset\LG)^{\wedge}
\]
上的拟凝聚层，这里用的是 $\LG$ 沿闭子集 $\LpG\subset\LG$ 的形式完备化。

\begin{remark}
考虑映射
\[
f\colon */\LG\to */(\LG)_{\dR},
\]
即从 $*/\LG$ 到其 de Rham 叠的映射。于是 $f_{1\sharp}\mathbf{1}$ 定义了 $*/(\LG)_{\dR}$ 上的一个范畴层，其底范畴（即沿 $*\to*/(\LG)_{\dR}$ 的拉回）是 $*/\LG^{\wedge}$，也就是 Kac--Moody 表示。此外，若取 $f_{1\sharp}\mathbf{1}$ 的球向量，即在范畴层面取 $\LpG$-不变量，我们恰好得到
\[
D_{\qc}(*/(\LpG\subset\LG)^{\wedge}).
\]
\end{remark}

\begin{remark}
几何局部 Langlands 对应应当是 $(\infty,2)$-范畴之间的一个等价。在 D-模设定下，它把 $*/(\LG)_{\dR}$ 上的范畴层，与去心圆盘上 $\Gd$-局部系统叠上的（$\IndCoh_{\Nilp}$ 版的）拟凝聚范畴层联系起来。在此等价之下，上一个注中的表示应当对应于\srcpage{51}局部 opers 叠上的层范畴。而取球向量则对应于过渡到无单值 opers。
\end{remark}

在 Arinkin、Beraldo、Campbell、Chen、Færgeman、Gaitsgory、Lin、Raskin 与 Rozenblyum（2024）中，在 Feigin 与 Frenkel（1992）、Beilinson 与 Drinfeld（1996）以及许多其他人的工作基础上，证明了如下定理（见该文定理 0.1.2）：

\begin{theorem}[Arinkin、Beraldo、Campbell、Chen、Færgeman、Gaitsgory、Lin、Raskin 与 Rozenblyum，2024]
存在交换图
\[
\begin{tikzcd}[column sep=large,row sep=large]
D(\Bun_G)_{\crit} \arrow[r,"\mathbb{L}_G"] & \IndCoh_{\Nilp}(\LocSys_{\Gd}) \\
\KL(G)_{\crit} \arrow[r,"\mathrm{FLE}_G"] \arrow[u,"\Loc_G"] & \IndCoh^{*}(\Op^{\monfree}_{\Gd})_{\Ran} \arrow[u,"\Poinc^{\mathrm{spec}}_{\Gd}"']
\end{tikzcd}
\]
其中下方水平箭头是一个等价，即「基本局部等价」（在临界水平上）。
\end{theorem}

这里，下标 $\crit$ 指的是「临界水平」，它是通常的 $\rho$-平移的一个体现。就 $D(\Bun_G)$ 而言，由于我们对 $\Th$ 的选取，它其实并不改变范畴。

函子 $\Poinc^{\mathrm{spec}}_{\Gd}$ 有一个与我们此前见过的 $\Poinc_{!}$ 的定义非常相似的描述。也就是说，利用（一般的）整体 opers 的概念，即只在有限多个点之外才有 oper 结构，可以引入局部 opers 的空间与整体 $\Gd$-局部系统的叠之间的一个对应。然后施加通常的拉-推函子。

至于 $\Loc_G$ 的定义，我们注意到：由覆叠 $*/\LG\to*/(\LG)_{\dR}$ 得到的 $*/(\LG)_{\dR}$ 的表示，其整体类比是由覆叠 $\Bun_G\to\Bun_{G,\dR}$ 得到的 $D(\Bun_G)=D_{\qc}(\Bun_{G,\dR})$ 中的对象。注意，这只不过给出 $\Bun_G$ 上那个恰为微分算子代数本身的 D-模；而 Kac--Moody 局部化接着利用 $\LG$ 的李代数的表示论，切出更为具体的 D-模。

\begin{remark}
粗略地说，这个图使人能够构造有趣的 Hecke 特征层。也就是说，给定一个带 oper 结构的 $\Gd$-局部系统，就能在右下角造出一个对象。在基本局部等价之下，它可以被搬到自守侧，然后施加 Kac--Moody 局部化 $\Loc_G$，便得到 $\Bun_G$ 上一个有趣的 D-模，而它实际上将是给定 $\Gd$-局部系统的 Hecke 特征层。
\end{remark}
\srcpage{52}
